% p1-chp01-introduction.tex

\chapter{Introduction}
\label{chp:introduction}

OctaC is a small, statically typed language for numerical and matrix
computation. Its purpose is to make numerical scripts, matrix operations,
and iterative computations easy to read and easy to write, without giving
up static typing.

This chapter introduces the language at a high level: the motivation behind
it, the languages it draws on, the features it offers, and the shape of its
syntax through a set of sample programs. The precise definition of its
tokens, and the lexer that recognizes them, follows in
Chapter~\ref{chp:lexer}.

\section{Motivation}
\label{sec:motivation}

Languages built for matrix work, such as Octave and MATLAB, make matrices
and vectors easy to write, but they are dynamically typed
\cite{octave-manual,matlab-docs}. A program that multiplies two matrices of
incompatible shapes is accepted as written and fails only when the operation
runs. Languages such as C and C++ check types before a program runs, but
treat matrices as library structures or nested arrays, and leave the
arithmetic to hand written loops \cite{kernighan1988c,stroustrup2013cpp}.

OctaC sits between the two. Matrices and vectors are built in types, and
their arithmetic is part of the language. This requires the semantic
analysis phase to track shape alongside type: operations on statically known
shapes are checked at compile time, while operations on shapes that depend
on a function parameter, a return value, or a loop built matrix are checked
at runtime.

\section{Design Influences}

OctaC is a hybrid rather than a subset of a single language. Each part of
its design comes from the language that handles it best:

\begin{itemize}
    \item \textbf{Octave and MATLAB} for matrix and vector arithmetic
    \cite{octave-manual,matlab-docs}.
    \item \textbf{C and C++} for block structure and statement syntax
    \cite{kernighan1988c,stroustrup2013cpp}.
    \item \textbf{Rust and Zig} for explicit sized primitive types
    \cite{rust-book,zig-reference}.
    \item \textbf{Python} for range based and \code{for...in} iteration, and
    for English word logical operators (\code{and}, \code{or}, and so on)
    \cite{python-reference}.
\end{itemize}

\section{Features}

\noindent\textbf{Native matrix and vector types.} Matrices and vectors are
built in types, not library structures. Operators are split between true
matrix arithmetic (\code{*} for multiplication, \code{'} for transpose) and
element wise arithmetic (\code{.+}, \code{.*}, and related operators), so
the operation being performed is visible at the point of use.

Beyond this, OctaC supports:

\begin{itemize}
    \item Explicit sized variable declarations (\code{i32}, \code{f32},
    \code{bool}, \code{matrix<i32>}, \code{matrix<f32>}, and related types)
    \item Block scoping with shadowing in nested blocks
    \item Scalar, matrix, and element wise operators
    \item Assignment, including compound forms (\code{+=}, \code{*=})
    \item Conditionals: \code{if} / \code{else if} / \code{else}, using
    English word logic (\code{and}, \code{or}, \code{not}, \code{is})
    \item Loops: range based \code{for}, \code{for...in}, \code{until},
    \code{do...until}
    \item Functions with typed parameters and an explicit return type
    \item Matrix and vector literals and indexing
\end{itemize}

\section{Sample Programs}

The following programs show the syntax of OctaC in use. Each one is a
complete program, followed by a walkthrough of the constructs it introduces.
Later programs build on what earlier ones explain, so they are best read in
order.

A few things are common to all of them. A program is made of functions, and
execution starts in the function named \code{main}, which returns an
\code{i32}. Statements end with a semicolon, and blocks are delimited by
braces, as in C. A comment starts with \code{//} and runs to the end of the
line. Output is written with \code{printf}, whose format string follows the
one in C \cite{kernighan1988c}: \code{\%d} prints an integer and \code{\%.2f}
prints a floating point number with two digits after the point. OctaC adds
\code{\%m}, which prints a matrix or a vector.

\subsection{Sum of Squares}

The first program adds up the squares of the numbers from 1 to 5. It shows the
basic building blocks of the language: a function, a typed variable, a loop,
and output.

\begin{lstlisting}
function main() -> i32 {
    i32 sum = 0;
    for (i32 i = 1...5) {
        sum += i * i;
    }
    printf("Sum: %d\n", sum);
    return 0;
}
\end{lstlisting}

The header \code{function main() -> i32} declares a function. The empty
parentheses are its parameter list, and the arrow introduces its return type,
so the type of the result is always written after the parameters. Inside the
body, \code{i32 sum = 0;} declares a variable. Every declaration in OctaC
follows the same order: the type first, then the name, then the initial value.

The loop \code{for (i32 i = 1...5)} is a range loop. It declares its own
counter, \code{i}, and counts it from 1 up to and including 5, using the range
operator \code{...}. Each pass adds \code{i * i} to the running total with the
compound assignment \code{+=}, so the program prints \code{Sum: 55}.

\subsection{Matrix Multiplication}

The second program shows the feature the language is built around, which is
that matrices are a native type rather than something a library provides.

\begin{lstlisting}
function main() -> i32 {
    matrix<i32> A = [[1, 2], [3, 4]];
    matrix<i32> B = [[5, 6], [7, 8]];
    matrix<i32> C = A * B;
    printf("A * B:\n%m\n", C);
    return 0;
}
\end{lstlisting}

The type \code{matrix<i32>} is a matrix whose elements are \code{i32}. A
matrix literal is written as a list of rows, and each row is itself a list in
brackets, so \code{[[1, 2], [3, 4]]} is a matrix with two rows and two
columns. Because \code{A} and \code{B} are initialized with literals at the
point of declaration, their shapes are known, and the declaration does not
need to state them.

The expression \code{A * B} is true matrix multiplication. Its value is
stored in \code{C}, and is the matrix \code{[[19, 22], [43, 50]]}. This
is different from element wise multiplication, which has an operator of its own.
The expression \code{A .* B} would multiply the entries position by position,
and would give the matrix \code{[[5, 12], [21, 32]]}. Because the shapes
of \code{A} and \code{B} are known, the compiler can check that they are
compatible for multiplication before the program runs. The format specifier
\code{\%m} prints the result.

\subsection{Matrix-Vector Product}

The third program combines the two native types by multiplying a $2 \times 3$
matrix by a vector of 3 elements ($A\mathbf{x}$), which produces a vector of
2 elements.

\begin{lstlisting}
function main() -> i32 {
    matrix<f32> A = [
        [1.0, 2.0, 3.0],
        [4.0, 5.0, 6.0]
    ];
    vector<f32> x = [2.0, 1.0, 0.0];

    // Linear transformation Ax -> 2x1 vector
    vector<f32> y = A * x;

    printf("Result y = Ax: %m\n", y);
    return 0;
}
\end{lstlisting}

Here \code{vector<f32>} is a vector of \code{f32} elements, written as a
single bracketed list. The matrix literal is spread over several lines, since
whitespace is not significant, and the comment marks what the multiplication
computes. The shape of the result follows the same rule as for two matrices:
it has as many rows as the left operand, and the multiplication is only valid
when the number of columns of \code{A} matches the number of elements of
\code{x}. Here both are 3, so \code{y} has 2 elements, namely
\code{[4.0, 13.0]}.

\subsection{Iterative Convergence}

The fourth program introduces conditions and the \code{until} loop. It halves a
value until it drops below a threshold, or until a step limit is reached.

\begin{lstlisting}
function main() -> i32 {
    f32 val = 100.0;
    i32 steps = 0;
    until (val < 1.0 or steps is 20) {
        val /= 2.0;
        steps += 1;
    }
    printf("Converged in %d steps.", steps);
    return 0;
}
\end{lstlisting}

An \code{until} loop runs as long as its condition is false, and stops as soon
as the condition becomes true. It is the reverse of the \code{while} loop of
C: the condition states when to stop rather than when to continue. The
condition here combines two tests. The comparison \code{val < 1.0} keeps its
symbolic form, while equality is written with the word \code{is}, and the two
tests are joined by \code{or}. The loop therefore ends when the value has
fallen below 1, or after 20 passes, whichever comes first. The second test
guards against a loop that would otherwise never end.

The body uses the compound assignments \code{/=} and \code{+=}. The value is
halved seven times before it falls below 1, so the program prints
\code{Converged in 7 steps.}

\subsection{Input Validation}

The fifth program reads a value from the user and keeps asking until it is
valid. It shows the \code{do...until} loop, which always runs at least once.

\begin{lstlisting}
function main() -> i32 {
    f32 val = 0.0;
    do {
        printf("Enter a nonzero val: ");
        val = input();
    } until (val is not 0.0);
    printf("Accepted: %.2f\n", val);
    return 0;
}
\end{lstlisting}

In a \code{do...until} loop the body comes first and the condition is checked
afterwards, so the prompt is always shown at least once. The function
\code{input()} reads a value from the user. The condition uses \code{is not},
the word form of inequality, so the loop repeats until the value entered is not
zero. The variable is given the initial value \code{0.0} so that it holds a
defined value before the first read.

\subsection{GPU-Accelerated Matrix Operations}

The last program shows two more parts of the type system: precision and
statically declared sizes.

\begin{lstlisting}
function main() -> i32 {
    // Statically sized matrix
    matrix<f16, 1024, 1024> A = identity();

    // can auto-infer size
    matrix<f16> B = A * 2.5;

    matrix<f16, 1024, 1024> C = A * B;

    return 0;
}
\end{lstlisting}

The element type \code{f16} is a half precision floating point number. In the
declaration of \code{A}, the dimensions are written as part of the type, so
this is a matrix of 1024 by 1024 elements. The matrix \code{B} is declared
without dimensions, because they can be inferred from the expression that
initializes it: multiplying \code{A} by the scalar \code{2.5} scales every
element and keeps the shape.

The last line multiplies two large matrices. Operations on matrices of this
size are meant to be offloaded to a GPU automatically when one is available,
without the program containing any GPU specific code.

\section{Looking Ahead}

The programs above show the surface of OctaC: typed declarations, native
matrices and vectors, two families of arithmetic operators, English word
conditions, and four kinds of loops. What they leave open is how a program
like these becomes something a machine can run. That is the job of the
compiler, which works in phases \cite{aho2006compilers}.

The first phase, lexical analysis, reads the program as a stream of
characters and splits it into tokens such as \code{for}, \code{i32}, \code{1},
and \code{...}. Chapter~\ref{chp:lexer} defines these tokens precisely and
shows how the lexer recognizes them. The later phases give the tokens a
structure and a meaning, and they are where the shape checking described in
Section~\ref{sec:motivation} takes place. They will be added to this book as
they are built.
